\documentclass[11pt]{article}
\usepackage[margin=1in]{geometry}
\usepackage{amsmath,amssymb,amsthm}
\usepackage{booktabs}
\usepackage{tabularx}
\usepackage{array}
\usepackage{enumitem}
\usepackage{microtype}
\usepackage{hyperref}
\usepackage[nameinlink,noabbrev]{cleveref}

\newcommand{\TV}{\mathrm{TV}}
\DeclareMathOperator{\KL}{D_{\mathrm{KL}}}
\newcommand{\MI}{I}
\newcommand{\PP}{\mathbb{P}}
\newcommand{\EE}{\mathbb{E}}
\newcommand{\one}{\mathbf{1}}

\newtheorem{definition}{Definition}
\newtheorem{remark}{Remark}
\newtheorem{proposition}{Proposition}
\newcolumntype{Y}{>{\raggedright\arraybackslash}X}

\title{Trace Interfaces for Anonymous DHTs:\\ Observation Tiers, Committee-Contact Traces, and Exposure Normal Forms}
\author{}
\date{Unified archive note v0.07 (2026-02-28; archive v139)}

\begin{document}
\maketitle

\begin{abstract}
Anonymous DHT lookups leak through \emph{interfaces}: what a passive observer (or the system itself via logs/certificates) can learn from the lookup transcript.
Different papers in this archive bound different projections of the transcript (a single committee label per epoch, a multiset of contacts, an ordered failure/timeout trace, \ldots).
This short note defines a compact, referee-facing \emph{normal form} for such projections.
It standardizes the committee-contact trace (CCT) object used across the Anonymous DHT profiling, state, certification, and auditing sub-series, and introduces a small ladder of observation tiers that lets papers stay non-redundant by \emph{naming} the interface they bound.
In receipts this normal form is referenced via a digest-bound \texttt{exposure\_nf\_id} (see Synthesis~12) and instantiated concretely in the worked example (Synthesis~17).
\end{abstract}

\paragraph{Scope.}
Define a compact normal form for transcript projections in Anonymous-DHT lookups (observation tiers and the committee-contact trace).

\paragraph{Imports.}
Baseline notation/model conventions are fixed by Synthesis~8.\cite{series:notation}
Receipt line items bind interface declarations via \texttt{exposure\_nf\_id}; see Synthesis~12.\cite{series:receiptitems}
This note defines the ENF/CCT objects that those fields point to.

\paragraph{Later-terse-paper citation posture.}
When a later short paper only needs to answer \emph{which transcript projection / tier / committee-contact-trace packet was declared for this Anonymous-DHT claim?}, it should stop at the minimal public ENF/CCT card below plus the tiny coarsening drill.
It should widen only when the live issue becomes the receipt tuple itself (Synthesis~12), the worked interlock across receipts and replay hooks (Synthesis~17), the horizon-accounting rulebook (Synthesis~19), or the downstream metric owners for Tier~0/1/2/3 claims (Anonymity~B, State~2, Boss Fight~A/B, and Congestion~1).\cite{series:receiptitems,series:workedexample,series:accountingrulebook,series:profiling,series:state2,series:bossfightA,series:bossfightB,series:congeq}

\begin{table}[t]
\centering
\small
\begin{tabularx}{0.97\linewidth}{@{}p{0.30\linewidth}Y@{}}
\toprule
\textbf{Public row} & \textbf{Pinned value}\\
\midrule
Illustrative ENF card id & \texttt{exposure\_nf\_id = sha256:7a1f\ldots}; the full digest belongs in the ABOM and receipts, but the card pins the exact canonical object being named rather than prose alone \\
Observation tier & Tier~1 contact multiset over public committee labels; this is strictly finer than a Tier~0 single-label claim and therefore cannot be silently downgraded in later prose \\
Alphabet / granularity & $\mathcal{A}$ is the release-bound committee partition with 16 public labels under \texttt{partition\_version = 2026.03}; changing the partition or coarsening depth mints a new ENF object \\
Repetition unit & One lookup epoch with reset at the lookup boundary; horizon claims must import Synthesis~19 rather than smuggling a different repetition unit into this card \\
Public conditioning & \texttt{cond} publishes the release id plus the public calendar / partition identifiers used to interpret the trace; if those public conditions change, the interface id changes too \\
CCT packet boundary & For this Tier~1 packet the CCT records \texttt{tier}, $\mathcal{A}$, \texttt{gran}, \texttt{unit}, and \texttt{cond}; background equalization law $Q$ and smoothing are left blank because they belong only when a downstream metric owner actually uses them \\
Owner routing & Synthesis~12 owns receipt line items, Synthesis~17 owns the fully interlocked worked example, and the downstream metric owners remain Anonymity~B / State~2 / Boss Fight / Congestion rather than this interface note \\
\bottomrule
\end{tabularx}
\caption{A minimal public ENF/CCT card. This is the smallest public answer later terse papers can quote directly when the live issue is which transcript projection was declared, rather than the downstream budget theorem, accountant, or release packet.}
\label{tab:minimal-enf-card}
\end{table}

This card is intentionally non-decorative: it pins one observation tier, one public alphabet/coarsening choice, one repetition unit, and one conditioning packet, without silently widening into the downstream metric owners or the broader receipt / replay stack.

\paragraph{One tiny coarsening drill.}
Suppose a later note wants to quote only the single winning committee label from the Tier~1 contact multiset declared above.
That later projection is a coarsening $g' = h \circ g$ of the original Tier~1 map, so the old budget may only be reused if the note mints a new \texttt{exposure\_nf\_id} for the Tier~0 label projection, keeps the same repetition unit and public conditioning, and declares the coarsening map explicitly.
If the repetition unit changes from ``one lookup epoch'' to ``one day of repeated lookups,'' the old card is no longer enough and the later note must reopen Synthesis~19 plus the relevant accountant owner.

\paragraph{Exports.}
(i) the observation-tier ladder, (ii) the CCT schema and fields, (iii) an ``exposure normal form'' for declaring which projection a bound applies to, and (iv) a stable ENF-ID scheme for binding certificates to that interface.

\paragraph{Artifact policy.}
Source-only; supporting artifacts are available on request. See \cite{series:artifactpolicy}.


\paragraph{Where this shows up downstream.}
Every receipt line item carries an \texttt{exposure\_nf\_id} to say which transcript projection its bound applies to.\cite{series:receiptitems}
Synthesis~17 provides a concrete instantiation (including ENF labels and digest IDs) so papers can stay non-redundant.\cite{series:workedexample}

\section{Why interfaces matter (and what this note owns)}
Anonymous lookup mechanisms typically protect payloads and queried keys, but what remains observable is a structured transcript: \emph{who was contacted, in what order, with what retries and termination rule}.
Attacks on DHT-based anonymity systems repeatedly exploit this transcript surface (passive range estimation, active selective failures, and long-horizon profiling) \cite{WangMittalBorisovCCS2010,MittalBorisov2012InfoLeaks,WangOctopus2012}.

\paragraph{Series ownership.}
This note owns \emph{interface vocabulary} for Anonymous DHT papers in this archive:
(i) the observation-tier ladder (\Cref{sec:tiers});
(ii) the committee-contact trace (CCT) object and its fields (\Cref{sec:cct}); and
(iii) the ``exposure normal form'' used to declare exactly which projection is being bounded.
All other papers should cite this note instead of re-defining tiers or the CCT schema.

\paragraph{Non-redundancy map.}
The profiling paper (Anonymity~B) budgets long-horizon leakage for Tier~0/1 destinations \cite{series:profiling}.
The Boss-Fight dial sheet compiles MaxL targets into knobs for Tier~2/3 transcripts \cite{series:bossfightA,series:bossfightB}.
The state sub-series (S1--S3) treats \emph{stateful} traces and what can safely be published \cite{series:state1,series:state2,series:state3}.
Certified menus and ABOM/OINL treat interfaces as declared certificate fields \cite{series:certmenus,series:abomoinl}.


\paragraph{Orthogonal: query privacy and modern contexts.}
The interface vocabulary here is agnostic to whether the destination key is hidden from routers (\emph{query privacy})
or merely unlinkable to an initiator.
Query-privacy mechanisms (e.g., OT-based approaches for robust DHTs) change the cryptographic payload surface but still leave a transcript projection that must be declared and budgeted. \cite{BackesQueryPrivacy2012}
Similarly, recent ``anonymous DHT'' proposals in blockchain settings inherit the same need to name the bounded projection and the repetition window. \cite{TsengNCDHT2024}


\section{Transcript-to-interface as a map}
A lookup episode produces a (possibly encrypted) \emph{raw transcript} $\tau$: the sequence of messages, endpoints, timing, failures, and termination.
An adversary rarely sees $\tau$ directly; instead it sees a projection
\[
\mathsf{Obs} = g(\tau) \in \mathcal{Y}
\]
for some (system- and threat-model-specific) map $g$.
All anonymity claims in this archive are statements about the induced distributions on $\mathcal{Y}$.

\begin{definition}[Exposure normal form]
A paper that claims to budget a lookup channel must declare:
\begin{enumerate}[leftmargin=*]
\item the observation alphabet $\mathcal{Y}$ (what is visible/logged);
\item the projection map family $g$ (which raw features are kept);
\item the repetition unit (epoch, lookup, round), i.e. how observations compose; and
\item any public conditioning (e.g., public calendars, committee partitions).
\end{enumerate}
We call this declaration an \emph{exposure normal form}.
\end{definition}

\begin{definition}[Exposure normal-form id (ENF-ID)]
An exposure normal form declaration is meant to be a \emph{claim field}, not prose.
To support portability and reuse across papers, deployments should compute a stable identifier:
\[
\mathsf{exposure\_nf\_id}=\mathrm{SHA256}(\mathrm{canon}(\mathsf{ENF}))\in\{0,1\}^{256},
\]
where $\mathrm{canon}(\cdot)$ is a canonical serialization of the tuple
$(\mathcal Y,g,\mathsf{unit},\mathsf{cond})$ (e.g., a JSON object with sorted keys, no whitespace, and explicit schema/version tags).
Publish $\mathsf{exposure\_nf\_id}$ as \texttt{sha256:<lowercase-hex>} (explicit algorithm tag). A short prefix may be used in text provided the full digest is in the ABOM and receipts (Anonymity~C).
\end{definition}

\begin{remark}[What ENF-ID buys]
ENF-ID makes ``the interface being bounded'' machine-checkable.
A notarized certificate or audit report that omits ENF-ID can silently shift to a different transcript projection.
Conversely, ENF-ID plus an explicit \emph{coarsening map} (the $h$ such that $g'=h\circ g$) lets certificates extend across interface changes without ambiguity.
\end{remark}


\begin{remark}[Coarsening is a first-class knob]
If $g'$ is a coarsening of $g$ (i.e. $g' = h\circ g$), then any distinguishing advantage, maximal leakage, or posterior-mass/PML bound for $g$ implies the same (or better) bound for $g'$ by data processing.
Thus ``what the interface reveals'' is itself an operator knob.
\end{remark}

\begin{proposition}[Portability requires interface identity or explicit coarsening]
\label{prop:interface-portability}
Consider a certified claim for exposure normal form $(\mathcal Y,g,\mathsf{unit},\mathsf{cond})$.
A downstream paper or deployment update may reuse that claim without re-certifying the metric only in two cases:
\begin{enumerate}[leftmargin=*]
\item the exposure normal form is unchanged up to renaming/version metadata; or
\item the new interface is a declared coarsening $g' = h\circ g$ under the same repetition unit and public conditioning.
\end{enumerate}
In Case~(2), every data-processing monotone endpoint used in this archive (distinguishing advantage, posterior-mass/PML, maximal leakage, TV) can only improve.
If the repetition unit, public conditioning, or projection map changes in any other way, a new interface id and a new claim object are required.
\end{proposition}

\begin{proof}[Proof sketch]
Case~(1) is tautological.
In Case~(2), the new observation is a measurable function of the old one, so data processing applies.
Changing the repetition unit or public conditioning changes the composed experiment itself, so the original accountant no longer certifies the updated claim.
\end{proof}

\paragraph{Operational consequence.}
Interface declarations should therefore carry a stable schema/version id and be treated like release-controlled claim fields rather than informal prose.
This is the interface-side analogue of contract-object versioning and release receipts.\cite{series:contractobjects,series:releaseproperty}

\section{Observation tiers for Anonymous DHTs}\label{sec:tiers}
We use a minimal tier ladder that is expressive enough to label the papers in this archive.
Tiers are not claims about adversary power in general; they are names for \emph{which projection of the transcript is being bounded}.

\begin{description}[leftmargin=*]
\item[Tier 0: destination label.] One label per epoch/lookup.
The observable is a single committee/shard label $O_t\in\mathcal{A}$.
This is the interface used by long-horizon committee profiling and equalization \cite{series:profiling}.

\item[Tier 1: contact set or multiset.] The observable is the set/multiset of committees contacted in an epoch,
$C_t\subseteq \mathcal{A}$ (or a multiset of size $w$).
This tier is the natural interface for MUCC-style contact-marginal equalization and conditional availability floors \cite{series:state2}.

\item[Tier 2: ordered trace with failures.] The observable is an ordered sequence of contact attempts with success/failure flags and (optionally) coarse timing buckets:
$T = ((a_1,f_1),\ldots,(a_m,f_m))$.
This tier captures selective-failure attacks and retry semantics \cite{WangMittalBorisovCCS2010}.

\item[Tier 3: packet/flow metadata.] The observable includes packet-level timing/size/flow features, which typically forces a different mechanism family (padding/scheduling) and a different witness (e.g., congestion witnesses) \cite{series:schedcompiler,series:congeq}.
\end{description}

\paragraph{Declaring a tier.}
A paper may bound multiple tiers if it explicitly says so.
Otherwise, statements should be read as tier-specific.
In particular, a Tier~0 equalization statement (single label per epoch) does \emph{not} imply a Tier~1 statement (contact multiset), and a Tier~1 statement does not imply a Tier~2 statement (ordered failures).

\section{The committee-contact trace object (CCT)}\label{sec:cct}
A recurring need in Anonymous DHT deployments is to bind an anonymity claim to a concrete, versioned interface.
We standardize a minimal object that can be embedded in ABOM/OINL and in auditing receipts.

\begin{definition}[Committee-contact trace (CCT) schema]
A CCT declaration is a tuple
\[
\mathsf{CCT} = (\mathsf{tier},\ \mathcal{A},\ \mathsf{gran},\ Q,\ \mathsf{smooth},\ \mathsf{unit},\ \mathsf{cond})
\]
with fields:
\begin{enumerate}[leftmargin=*]
\item $\mathsf{tier}\in\{0,1,2,3\}$: observation tier (\Cref{sec:tiers}).
\item $\mathcal{A}$: the public label alphabet (committee/shard labels after coarsening).
\item $\mathsf{gran}$: granularity/coarsening parameters (e.g., prefix depth, committee partition version).
\item $Q$: the public reference/background distribution when used (Tier~0/1 equalization), including how it is produced (specified vs measured).
\item $\mathsf{smooth}$: smoothing/trimming rule that enforces a floor (e.g., $Q_{\min}>0$ or an $(\varepsilon,\delta)$ tail format).
\item $\mathsf{unit}$: repetition unit and checkpoint rule (epoch length, lookup boundary, reset conditions).
\item $\mathsf{cond}$: any declared public conditioning (calendar ids, public randomness beacons, release ids).
\end{enumerate}
\end{definition}

\begin{remark}[Minimality]
This schema is intentionally small.
Anything protocol-specific (e.g., Kademlia $\alpha/\beta/k$ parameters, stop-time kernels, retry timers) belongs in the mechanism paper and is cited as such.
The CCT object exists to pin down \emph{what was bounded} so non-redundant papers can safely compose.
\end{remark}

\section{Interface-to-budget bridges (one-paragraph pointers)}
This note does not re-derive budgets; it only names the bridges so other papers can cite them.

\paragraph{Tier 0.}
Equalization and long-horizon profiling budgets for $O_t\in\mathcal{A}$ are in Anonymity~B \cite{series:profiling}.
Worst-case posterior-mass/PML caps require smoothing (a floor on $Q$); typical leakage can be expressed as mutual information scaling like $\Theta(T/w^2)$ under blanket mixing.

\paragraph{Tier 1.}
Contact-set privacy notions (MUCC and beyond) and the associated availability floors are in State~S2 \cite{series:state2}.
When a system logs a multiset (or an ordered set), it must be budgeted at this tier even if it later publishes only a single label.

\paragraph{Tier 2.}
Retry/abort semantics and selective-failure amplification are handled as repeated-use witnesses in the Boss-Fight series \cite{series:bossfightA,series:bossfightB} and as stateful witnesses in State~S1 \cite{series:state1}.

\paragraph{Tier 3.}
Timing/queueing/capacity witnesses are handled by the scheduling/padding backbone papers and the congestion series \cite{series:schedcompiler,series:congeq}.

\section*{Acknowledgement of limits}
The tier ladder is not a complete taxonomy: real deployments combine tiers (e.g., packet metadata plus committee labels).
The intent is pragmatic: a short, stable vocabulary that prevents papers from silently talking past each other.

\begin{thebibliography}{99}\small
\bibitem{series:artifactpolicy}
Anonymous.
\newblock Artifact and Disclosure Policy for the One-True-Archive (Source-Only Public Release).
\newblock Series synthesis note (Synthesis~6), 2026. (This archive.)

\bibitem{WangMittalBorisovCCS2010}
Q.~Wang, P.~Mittal, and N.~Borisov.
\newblock In search of an anonymous and secure lookup: attacks on structured peer-to-peer anonymous communication systems.
\newblock In \emph{ACM CCS}, 2010.

\bibitem{MittalBorisov2012InfoLeaks}
P.~Mittal and N.~Borisov.
\newblock Information leaks in structured peer-to-peer anonymous communication systems.
\newblock \emph{ACM TISSEC}, 15(1), 2012.

\bibitem{WangOctopus2012}
Q.~Wang and N.~Borisov.
\newblock Octopus: a secure and anonymous DHT lookup.
\newblock In \emph{IEEE ICDCS}, 2012.

\bibitem{series:profiling}
Anonymous.
\newblock Anonymous DHT Profiling and the Equalization Contract.
\newblock Anonymity series Paper~B, 2026. (This archive.)

\bibitem{series:abomoinl}
Anonymous.
\newblock ABOM and OINL for Anonymity Claims: Proof-Carrying Odds--Inflation Budgets Without Publishing Full Artifacts.
\newblock Anonymity series Paper~C, 2026. (This archive.)

\bibitem{series:bossfightA}
Anonymous.
\newblock Boss Fight Budgets: Maximal Leakage as a Repeated-Use Anonymity Contract.
\newblock Boss Fight series Paper~A, 2026. (This archive.)

\bibitem{series:bossfightB}
Anonymous.
\newblock AnonDHT in the Boss Fight regime: compiling a MaxL budget into lookup knobs.
\newblock Boss Fight series Paper~B, 2026. (This archive.)

\bibitem{series:certmenus}
Anonymous.
\newblock Certified menus for anonymous DHT lookups.
\newblock Certified series Paper~A, 2026. (This archive.)

\bibitem{series:state1}
Anonymous.
\newblock State-dependent anonymity: selection and state evolution as an attack surface.
\newblock AnonDHT state series Paper~1, 2026. (This archive.)

\bibitem{series:state2}
Anonymous.
\newblock Marginally uniform committee contact in anonymous DHT lookups: availability floors, perfect-leakage counterexamples, and abort-safe equalization.
\newblock AnonDHT state series Paper~2, 2026. (This archive.)

\bibitem{series:state3}
Anonymous.
\newblock Closed-view auditing for stateful anonymity: odometers, alarms, and machine-checkable CPPCs.
\newblock AnonDHT state series Paper~3, 2026. (This archive.)

\bibitem{series:schedcompiler}
Anonymous.
\newblock Scheduling and compiler techniques for metadata-hiding overlay lookups.
\newblock 2026. Wiki: \texttt{[[2026.01.22 - Mathematics: Scheduling and Compiler Techniques for Metadata-Hiding Overlay Lookups]]}.

\bibitem{series:congeq}
Anonymous.
\newblock Congestion-EQ: congestion witnesses and accounting.
\newblock Congestion series Paper~1, 2026. (This archive.)

\bibitem{series:contractobjects}
Anonymous.
\newblock Contract Objects for Anonymous-DHT Privacy Budgets and Claims.
\newblock Series synthesis note (Synthesis~0), 2026. (This archive.)

\bibitem{series:releaseproperty}
Anonymous.
\newblock Release-Property Ledgers and Receipts for Anonymous-DHT Claims.
\newblock Release and destination series Paper~A, 2026. (This archive.)

\bibitem{BackesQueryPrivacy2012}
M.~Backes, I.~Goldberg, A.~Kate, and T.~Toft.
\newblock Adding query privacy to robust DHTs.
\newblock In \emph{Proc.\ ACM ASIACCS}, 2012.

\bibitem{TsengNCDHT2024}
L.~Tseng, L.~Ambarapu, and M.~Aloqaily.
\newblock NC-DHT: a robust and anonymous DHT for blockchain systems.
\newblock In \emph{Proc.\ IEEE International Conference on Blockchain Computing and Applications (BCCA)}, 2024.


\bibitem{series:notation}
Anonymous.
\newblock Notation and Minimal Model for the One-True-Archive.
\newblock Series synthesis note (Synthesis~8), The-One-True-Archive (2026).

\bibitem{series:receiptitems}
Anonymous.
\newblock Receipt Line Items for Anonymity Claims: A Minimal Tuple Schema for Published Budgets.
\newblock Series synthesis note (Synthesis~12), The-One-True-Archive (2026).

\bibitem{series:workedexample}
Anonymous.
\newblock Worked Example: Instantiating a Tiered Reader-Privacy Receipt.
\newblock Series synthesis note (Synthesis~17), The-One-True-Archive (2026).

\bibitem{series:accountingrulebook}
Anonymous.
\newblock Receipt Accounting and Horizon Composition for Anonymous-DHT Budgets.
\newblock Series synthesis note (Synthesis~19), The-One-True-Archive (2026).

\end{thebibliography}

\end{document}
